\subsection[Challenge 57: LIGO with modified mirrors for ultralight vector dark matter]{Challenge 57: LIGO with modified mirrors for ultralight vector dark matter}
\textbf{Problem ID:} \texttt{Challenge\_57\_main}\\
\textbf{Primary inferred discipline:} Gravitation, Cosmology \& Astrophysics\\
\textbf{Secondary inferred disciplines:} High Energy \& Nuclear Physics\\
\textbf{Python implementation:} \texttt{challenges/57/answer.py}

\subsubsection*{Problem}
\paragraph*{Problem setup}
Consider a vector field $A^\mu$ coupled to the standard model $B-L$ current $J^\mu_{B-L}$ through $\mathcal{L}\supset -\epsilon_{B-L}eJ^\mu_{B-L} A_\mu$. If the vector field makes up all the dark matter, the local dark matter density indicates that the field could exert a measurable force on mirrors in laser interferometers, with sufficiently large coupling. Consider LIGO, with a strain sensitivity of $3\times 10^{-24}\, \text{Hz}^{-1/2}$ at frequency $250\,\text{Hz}$.

\paragraph*{Main problem}

Take the dark charge $Q_D$ to mass $M$ radio to be $\frac{Q_D}{M}\sim \frac{0.5}{m_n}$ for the inner mirrors, where $m_n$ is the mass of a neutron. We introduce a doped outer mirror, resulting in a differential forced between the inner and outer mirror. We assume the outer mirror gains an additiona $\delta \left(\frac{Q_D}{M}\right)\sim \frac{\delta q}{m_n}$, where $m_n$ is the mass of a neutron. With an observation time of $13$ years and SNR of $1$, what is the smallest $\epsilon_{B-L}$ it can probe at $250\,\text{Hz}$? Provide answers for scenarios where $\delta q=\{0.074, 6\times 10^{-3}, 5\times 10^{-4}\}$.


\paragraph*{Reference-model assumptions}
Use $\rho=0.4\,\mathrm{GeV\,cm^{-3}}$, dark-matter dispersion and streaming
speed 230 km/s, arm length $L=3994.5$ m, fused-silica inner mirrors, and the
specified inner-mirror charge ratio $Q_D/M=0.5/m_n$. Model suspended mirrors
as free masses. Keep the undoped finite light-travel-time response and the
spatial-gradient contribution as well as the added doping response.
Average over independent isotropic field polarization and propagation
orientations, with polarization randomized between coherence patches.
Use a continuous 13-Julian-year record, the stated one-sided strain-noise
spectral density, and the semicoherent power convention
\[
\tau=\frac{1}{f(v/c)^2},\qquad
T_{\rm eff}=\sqrt{\tau T},\qquad
\langle h^2\rangle_{\rm threshold}=S_h/T_{\rm eff}.
\]
Thus SNR one means equality of signal and noise in this power statistic;
do not assume phase coherence over all 13 years or an optimal fixed
polarization. Report the dimensionless coupling using SI electromagnetic
normalization, with the elementary charge $e$ as written in the interaction.
